\documentclass[11pt,a4paper]{article}

% --- Common Packages across layouts --- %
\usepackage{microtype}
\usepackage{graphicx}
\usepackage[hidelinks]{hyperref}
\usepackage{xcolor}
\usepackage{todonotes}
\usepackage{dirtytalk}
\usepackage{siunitx}
\usepackage{xspace}
\usepackage{lipsum} % For testing

%\usepackage{allrunes} package allrunes to add, but currently unstable due to trying to redefine \swdefault and \dot



% Article and report specific packages %

% Visual style packages

\usepackage[margin = 2.5cm]{geometry}
\usepackage{wrapfig}
\usepackage{caption}
\usepackage{subcaption}
\usepackage[nocheck]{fancyhdr}
\usepackage{titlesec}
\usepackage{blindtext}
\usepackage{enumitem}
\setitemize{itemsep=2pt, topsep=4pt, parsep=4pt, partopsep=6pt}

% Maths / Science packages
\usepackage{amsfonts, amssymb, amsmath, amsthm}
\usepackage{bm}
\usepackage{physics}
\usepackage{tikz}
\usepackage{circuitikz}
\usepackage{cancel}
\usepackage{mathrsfs} % New font family from \mathscr{}

% Other utility packages
\usepackage{multicol}
\usepackage{afterpage}
\usepackage[backend=biber, style=numeric]{biblatex}
\addbibresource{references.bib}



% My (John's) personal utility packages
\usepackage{packages/johns-package-setup}
\usepackage{packages/johns-preferences}


% ============================================ %
%    ARTICLE / SHORT REPORT TEMPLATE         %
% ============================================ %

% Article/short report template specific tweaks -- consider turning into a new (conditional) package
\numberwithin{equation}{section}

% Setting up headers and footers for fancyhf
\fancyhf{}
\fancyhead[R]{\today}
\fancyhead[L]{Project Overview}
\fancyfoot[C]{\thepage}
\headsep = 15pt
\voffset = 10pt

% Title and introduction to the report
\title{Project Overview}
\author{John Perry\thanks{Churchill College}}
\date{\today}

\begin{document}

\pagestyle{empty}
\newgeometry{margin = 2cm}
\pagenumbering{arabic}

\maketitle

% Set headers and footers to exist past the title page
% \afterpage{\thispagestyle{fancy}} % Was useful when not dealing with multicol
\pagestyle{fancy}

\begin{abstract}
    % Give the abstract of the document with the reason it exists, the significance of its contribution to the project and the conclusions it provides.
    Document containing the overview of the plans and goals of the Custom Split Keyboard project. First, the different kinds of property desired from the keyboard are listed with a brief summary of each, then following that a more in-depth analysis of each property is given.
\end{abstract}

\section{Introduction}
% Introduce the problem, and contextualise it and the objectives of the experiment/thing being reported on
I currently feel like the keyboards I have don't fully meet my needs practically. Key things include:
\begin{itemize}
    \item More optimised/better battery life.
    \item More portable so it can be taken round different places more easily.
    \item More comfortable for the work I want to do with it.
\end{itemize}

\subsection{Objectives}

The different objectives for the finished keyboard include:
\begin{itemize}
    \item \textbf{Ortholinear}: Partly that I want to try it out, but also it is meant to be very comfy.
    \item \textbf{Custom layout}: I want to have a comfy layout that is much nicer for coding but also not too dissimilar from QWERTY because I realise that simply fully going away from the norm will be a massive pain to learn again.\footnote{Consider how much more I can use my thumbs for, as just on the spacebar is a waste!}
    \begin{itemize}
        \item Additionally, it would be good to have a \textit{flexible} layout as well. I want to be able to change my opinions on the locations of different keys and experiment.
        \item \SI{80}{\percent}:\footnote{Or smaller? It should be investigated if I am comfortable with using layers \textit{if} they are comfortable for me.} I want a keyboard that is not too large.
    \end{itemize}
    \item \textbf{Split layout}: For ergonomics I want to try out a split layout, but I also want to be able to connect the two sides in a robust way to make a single keyboard (that also looks sleek) if wanted.\footnote{If I can connect the two sides with pogo pins that would be quite cool.}
    \item \textbf{Low profile}: Currently I don't have any low profile keyboard, and it would be cool to have something that is more portable and practical to use.
    \item \textbf{Top-of-the-line switches}: I might want to be able to use the keyboard for gaming, so top-tier switches would be preferable.
    \item \textbf{Compact and portable}: I want to be able to pack and move the keyboard easily, to actually make it feel good to use and that I would want to take it with me.
    \item \textbf{Good battery}: Since the battery is a point where many keyboards seem to make a compromise and cut corners, I want the battery life to be exceptional even in a small chassis.
    \item \textbf{Nice feeling materials}: I want the keyboard to \textit{feel} premium. That means ceramic-like keycaps, metal (or high-quality plastic) frame, and nice rubber feet to make it feel sleek and premium.
    \item \textbf{Good sound/acoustics}: I want it to feel nice to use, so it should have the capability to sound good---using nice materials and including space for sound dampening/adjustment.
\end{itemize}

\section{Custom layout}

\begin{itemize}
    \item Use my thumbs for more! What would be good to have on my thumbs would be commonly used coding symbols/buttons to press. For example:
    \begin{itemize}
        \item \textit{Enter} key: When writing code, it is very inconvenient to move fingers to press enter. It is also very on-theme.
        \item \textit{Brackets/braces}: Different brackets and braces (e.g. [ ], \{ \}, ( )) would be good to have on thumbs. Particularly having the opening ones on my left and closing on my right would feel quite nice I'd imagine.
    \end{itemize}
    \item Relocating \say{home} and \say{end} might actually work well in a workflow if I get in the right patterns.
    \item As well as pure coding, it would be nice to have 
\end{itemize}

\subsection{Ideas for the layout}

\begin{itemize}
    \item Decide if I want a spacebar on both thumbs. Currently the spacebar is only hit with my \textit{right thumb}, therefore having such a large area used by it is a waste.
    \item I definitely want shift keys on both sides, as capitalising is now an ingrained habit.
    \begin{itemize}
        \item Consider using the shift keys like a \say{\textit{function}} key, as the shift key already works analogous to one. If I don't want dedicated keys for the function row, consider how the symbol keys might be moved if the function keys take the place of the \say{shift} layer on the number keys.
    \end{itemize}
    \item New thumb items, as currently they are underused for a particularly dexterous finger:
    \begin{itemize}
        \item Enter key, as during coding it is inconvenient to move there when writing code.
        \item Open/close brackets/braces: Depending on how I want to use layers, (and if I have 2 layers with a \say{shift} layer as well as a \say{function} layer) I could get these out of the way with one key on either side.\footnote{For ( ), [ ], \{ \}.}
    \end{itemize}
    \item Layers: Shift keys already provide another layer, and it makes sense to have one on either side. Consider adding a \say{function} key for a layer on either side too.\footnote{This could be on the underused thumbs, for example.}
\end{itemize}

\subsection{\texorpdfstring{\SI{80}{\percent}}{80 Percent}}

\SI{80}{\percent} should be considered more of a rough goal, I want a keyboard that is flexible enough that it doesn't take up too much space. \SI{80}{\percent} has been seen to be a good balance in the past, but it must be considered which keys I want immediate access to and which I am OK to have on layers, as this informs the necessary size for the keyboard.

\subsection{Wrist-rests}

Having wrist-rests would add to the premium feel of the keyboard. Some initial comments on wrist-rests are:
\begin{itemize}
    \item I want them to be integrated into the base design, as carrying around a separate wrist-rest feels clunky and I don't like it.
    \begin{itemize}
        \item Using magnets may offset this feeling slightly, which may also be considered.
    \end{itemize}
    \item Depending on how low-profile I manage to get the keyboard, it could be that I don't need a wrist-rest.\footnote{Such as is the case for the \href{https://www.apple.com/uk/shop/product/mxcl3b/a/magic-keyboard-usb-c-british-english}{Apple Magic Keyboard}.}
    \item The wrist-rest would have to be able to (at the very least) split in half, as it would need to both function with both sections together and apart.\footnote{For inspiration, the split \href{https://nuphy.com/products/wrist-rest-for-node-series-low-profile?variant=43916448137325}{Nuphy Wrist-rests} could be a source of inspiration.}
\end{itemize}

\section{Switches and Keycaps}

To meet as many of the desired characteristics laid out in the project objectives, some compromises may need to be made. For example, to achieve the low-profile nature of the keyboard, low-profile switches are preferable however for top-of-the-line switches the best that exists for gaming right now are hall-effect switches, which may not exist in a low-profile format. Thus, a compromise is needed.

\subsection{Keycaps}

Keycaps as an object can be fairly flexibly produced to an exact spec using 3D-printing techniques. More specifically, resin 3D printing could produce high-tolerance custom keycaps with mechanical properties that make them \say{feel} nice to use.\footnote{With properties such as density, and acoustic properties.}

\subsubsection{First ideas}

Some first ideas for producing nice-feeling keycaps are:
\begin{itemize}
    \item Resin can cause skin irritation long-term, so they should be coated in some way to prevent this.
    \item If coating the keycaps, it could be possible to coat in a substance that imitates the feeling and look of ceramic keycaps.\footnote{This may be in the form of a coloured addition, or a clear-coat.}
    \item Additionally, it could be possible to fill the letter divots with another colour which would give a really professional look.
    \item Curing the keycaps without washing can give a smoother look on the surface, at the cost of a little extra inconvenience.
    \item Some further discussion of resin keycaps can be found in \href{https://www.reddit.com/r/resinprinting/comments/thtnry/resin_printed_keycaps_for_my_keyboard/}{this Reddit thread}.
\end{itemize}

\section{Project Organisation}

Project organisation will be slightly altered as the project goes on, as it coincides with setting up my best practice workflows. These changes will be listed clearly here, to help with version control and seeing these documents in the future.

\subsection{CAD software of choice}

As part of making a switch to Linux during this project, the issue of CAD software came up---as there isn't a good \SOLIDWORKS alternative on Linux---so to respond to this, \ONSHAPE was used as it only relies on the browser and retains some of the more advanced features that exist in \SOLIDWORKS.

Because of this, however, in \texttt{/CAD/} there exist orphaned \SOLIDWORKS files which have then been interpreted further in \ONSHAPE, using the \ONSHAPE built-in version control system.

\subsubsection{What this means}

To ensure that this project remains in an open-source format, where someone can just copy and paste the instructions contained within to reproduce the project, upon project completion\footnote{Or on a given prototype version} the \texttt{stl} files from \ONSHAPE must be downloaded zipped to be uploaded to \GITHUB for publishing.

\end{document}